% === 基本設定 ===
\usepackage{graphicx, xcolor}       % 画像の挿入・文字色
\usepackage{tikz}

% === フォント設定 ===
\usepackage{unicode-math}
\usepackage{luatexja-fontspec}

% --- 欧文・数式設定 ---
\setmainfont{TeX Gyre Termes}
\setsansfont{TeX Gyre Heros}[Scale=MatchLowercase]
\setmathfont{TeX Gyre Termes Math}

% --- 和文設定 ---
\usepackage[haranoaji, deluxe]{luatexja-preset}

% === レイアウト・装飾 ===
\usepackage{geometry}
\geometry{
  top=25mm,
  bottom=25mm,
  left=20mm,
  right=20mm
}
\pagestyle{plain}                   % フッター中央にページ数
\usepackage{float}
\usepackage{multirow}               % 列の結合
\usepackage{booktabs}               % きれいな表の罫線
\usepackage{tabularx}               % 表の幅自動調整
\usepackage{enumitem}               % 箇条書きの高度な制御
\usepackage[normalem]{ulem}         % 強調を上書きしない設定
\usepackage{tcolorbox}              % 文章囲み
\tcbuselibrary{breakable}           % tcolorboxのページ跨ぎ
\usepackage[mathlines]{lineno}

% === 数式・物理 ===
\usepackage{amsmath}                % 数学記号・環境
\usepackage{physics}                % 物理系記法
\usepackage{siunitx}                % SI単位系の統一的表記
\AtBeginDocument{\RenewCommandCopy\qty\SI} % physicsとの競合回避

% === 図表・式の番号設定 ===
\usepackage[compatibility=false, font=small, labelfont=bf]{caption}
\usepackage{subcaption}             % 図の並列配置
\usepackage{listings}               % プログラムコードの貼り付け
\lstset{
  basicstyle=\ttfamily\small,       % 基本のフォント
  commentstyle=\color{green!60!black},  % コメントの色
  keywordstyle=\color{blue},        % キーワードの色（int, forなど）
  stringstyle=\color{purple},       % 文字列の色
  frame=single,                     % 枠で囲む
  breaklines=true,                  % 長い行を折り返す
  numbers=left,                     % 行番号を左に
  numberstyle=\tiny\color{gray},    % 行番号のスタイル
  stepnumber=1,                     % 行番号を1行ごとに表示
  columns=fixed,                    % 文字幅を等間隔に
  basewidth=0.5em,
  showstringspaces=false,           % 文字列中のスペース記号を表示しない
}
\numberwithin{equation}{section}    % 数式番号を節ごとに
\counterwithin{table}{section}      % 表番号を節ごとに
\counterwithin{figure}{section}     % 図番号を節ごとに

% === リンク・参照 ===
\usepackage{hyperref}               % ハイパーリンク
\usepackage{url}                    % URLを綺麗に表示
\usepackage{cleveref}               % 参照

% --- cleverefの日本語設定 ---
\crefname{equation}{式}{式}
\crefname{figure}{図}{図}
\crefname{table}{表}{表}
\crefname{page}{ページ}{ページ}
\crefname{section}{節}{節}
\crefname{appendix}{付録}{付録}
\creflabelformat{equation}{#2#1#3}

% === その他 ===
\usepackage{pxrubrica}              % ルビ振り

% === タイトルデザイン ===
\makeatletter
\renewcommand{\@maketitle}{
  \newpage
  \null
  \begin{center}
    \sffamily\gtfamily              % ゴシック・サンセリフ
    \@title                         % タイトル出力
  \end{center}
  \begin{flushright}
    \sffamily\gtfamily              % ゴシック・サンセリフ
    \large                          % 著者名の文字サイズ
    \@author                        % 著者名出力
  \end{flushright}
}
\makeatother